\section{Training-free selective explanation: protocol and novelty}
The next study was explicitly authorized after the completed abstention pilot. Protocol, runner and descriptive endpoints were frozen and pushed at commit \texttt{7778d09} before inference. The three conditions reuse all 24 previously consulted UCF inputs. This is a paired engineering diagnostic on five-second windows, not a full original-VERA evaluation or untouched test. No new human annotations are collected. The existing UCA captions are partial references, sometimes including nonvisual inferences.

\begin{center}\begin{tabularx}{\linewidth}{lXX}\toprule
Condition & Binary anomaly decision & Explanation intervention\\\midrule
A: Original & Original 0/1 instruction & Original questions and free-text answer\\
B: Manual & Best 0/1 decision & Observable claims plus a separate list of unresolved details\\
C: Cited & Same as B & Same schema, with supporting frame-slot citations requested for each observation\\\bottomrule
\end{tabularx}\end{center}

The same frozen InternVL2-8B weights and original five guiding questions are shared by all conditions. B and C remove permission to make assumptions and use the same JSON explanation format. Thus A--B is an instruction/format bundle, while B--C isolates the additional citation request more closely. Each condition generates its own binary decision; we measure preservation instead of forcing agreement. The new study does not allow decision 2: uncertainty is attached to explanation details, not to the final anomaly label. One shared model is loaded on Ada5 and calls run sequentially. No optimizer, gradients, LoRA, verbalized question updates, retrieval or external judge is used.

The bounded allocation is 72 calls, bringing the two pilots to 120 of the 192-call ceiling. The original baseline is rerun and compared with its previous responses. Exact prompts, input and dependency hashes, raw responses, strict binary parses and explanation-schema failures are retained. Invalid explanation JSON counts as a failure and is not silently repaired or retried.

\textbf{Novelty assessment.} The narrow method claim is not established. GroundedVQA already associates claims with video timestamps and flags unverified claims after evidence checks (\url{https://arxiv.org/html/2608.15574v1}, Section 3). TRACE also generates evidence-linked claims (\url{https://arxiv.org/html/2605.16740v1}, Section 3). Anomaly-specific visual verification and conservative score preservation appear in \url{https://arxiv.org/html/2610.16324v1}. Our single-call prompts provide less verification than these systems; adding citations alone is not new. The present run can test operational coupling between detection and explanation, but cannot establish a publishable method.

\textbf{What existing annotations can measure.} We report positive-reference content-word recall against UCA captions for the twelve annotation-included windows, with the word sets and reference hashes frozen before inference. For A, words come from its answer text; for B/C, only observation claims. This is lexical overlap, confounded by negation, synonyms, output length and formats. It is not semantic support or hallucination accuracy. Gap claims are never labelled false due to annotation omission. Citation indices in 1--8 establish syntax only. Timestamp overlap is deliberately not scored: every annotation-included frame lies in the selected interval, making that score structurally trivial.

\section{How original VERA learns and makes decisions}
Original VERA's ``learning'' can optimize text while leaving neural-network weights frozen. In a learner call, the model sees frames, its binary instruction and the current guiding questions and generates predictions. In a separate optimizer call, the same checkpoint receives the current questions, predictions and target labels and proposes revised question text. The program stores this text for future calls; the model does not rewrite its own weights. A single GPU can serve both roles sequentially. Validation selects the question set used later. See \url{https://vera-framework.github.io/} and the pinned original optimizer instruction.

At inference, the instruction, saved questions and sampled frames are sent with each model call. The model generates answer text and a label, and code parses the final Output field. The meaning of 0/1, or the modified 0/1/2 contract in the first pilot, is supplied in that prompt; it is not an additional trained classifier. Original full-video VERA subsequently uses scene-context fusion and temporal postprocessing, which neither small pilot runs. No verbalized-learning or neural-network training took place in either pilot.

\subsection{Preserved first attempt and frozen repair}
The first 72-call selective run completed in 190.54 seconds. All 24 original responses exactly matched the previous baseline. Manual and cited instructions both retained only five of twelve original alerts (four of seven on annotation-included inputs). These are changes in decisions, not measured accuracy losses or gains.

A prompt-template defect prevented meaningful explanation interpretation: the requested JSON example contained literal text ``a concrete visible observation,'' which the model copied into four of its five schema-valid cited observations. The manual condition had eighteen schema-valid responses but no populated observations in those valid objects; the cited condition had twenty-two schema-valid responses. One unresolved item also copied the template wording. Syntax compliance alone would have overstated usefulness. These original records and their original metrics remain preserved.

Amendment01 was frozen and pushed at \texttt{f355121} before a single rerun. It removes example values, asks for one to three actual visible observations even in normal scenes, and rejects empty observation lists or copied example strings. Requiring a nonempty observation is an additional instruction change; cross-run differences cannot be attributed to punctuation alone. No questions, weights, input selection or reference thresholds change. This is an output-informed manual repair on already consulted diagnostic sources, not blinded evaluation or a training result. The extra seventy-two calls bring the complete allocation to 192; no further repairs or learned extension are permitted within this budget.
